\chapter{Discussion}
\label{ch:discussion}

\section{Interpretation of the findings}
\label{sec:interpretation}

{Write once the results exist. Answer each research question directly, then
explain what the findings mean for each audience: certification bodies, vendors,
buyers, and researchers who assess third-party cryptographic dependencies.}

A long certification lag is not automatically a security failure: a vendor may
backport fixes without changing the version string, and re-validating a module has
a real cost. The lag should therefore be read as a measure of the distance between
the certified and upstream states rather than as a direct measure of risk.

\section{Threats to validity}
\label{sec:threats}

\paragraph{Incomplete disclosure.} The largest threat is that the absence of a mention
does not imply the absence of a dependency, because vendors may redact or generalize
their dependency descriptions. All prevalence figures are therefore lower bounds on
actual use, and disclosure practices may differ between vendors, schemes, and
periods. An apparent trend may reflect changed disclosure rather than changed use.

\paragraph{A mention is not an evaluated dependency.} A library may appear in a document
without being evaluated, or without being within the certified boundary. The evidence
levels in Table~\ref{tab:evidence} mitigate this threat but do not remove it, since the
levels are based on what the documents state.

\paragraph{Extraction errors.} Text extraction from PDF documents can fail or corrupt
content, and both regular expressions and language models produce errors. Measuring
precision and recall on an annotated sample bounds these errors. The sample itself is
subject to annotator judgment and to sampling error.

\paragraph{Version semantics.} Version strings can hide backported patches, vendor builds,
or a validated provider that differs from the library using it, which affects the
lag defined in Equation~\ref{eq:lag}. The lag can also only be computed for certificates
with a stated version, and these may not represent all certificates.

\paragraph{Dataset dependence.} The analysis inherits any errors and coverage limits of
sec-certs. Pinning a snapshot makes the study reproducible, but the findings describe the
state of the data at that date.

\paragraph{Non-independent observations.} One vendor may hold many closely related
certificates, so raw certificate counts can overweight large product families. Results
should therefore also be reported per vendor or product family where this changes the
conclusion.

\section{Limitations}
\label{sec:limitations}

The study relies only on public documents. It does not inspect binaries or source code,
so no statement about a product's dependency is confirmed against the product itself.
It also does not assess the implementation quality of any library or whether a
particular product is vulnerable.
